\section{Variables}

In Go (Golang), a variable is a named storage location used to hold a value
that can be accessed and changed while a program is running. Every variable
has a data type, which determines what kind of value it can store and how that
value is handled by the program.

Variables can be declared using the \texttt{var} keyword. For example:

\begin{lstlisting}
var age int
age = 25
\end{lstlisting}

This declares a variable called \texttt{age} with the \texttt{int} data type
and assigns it the value \texttt{25}. A variable can also be declared and
initialised at the same time:

\begin{lstlisting}
var age int = 25
\end{lstlisting}

Go can infer the data type from the assigned value, so the type can often be
omitted:

\begin{lstlisting}
var age = 25
\end{lstlisting}

Inside a function, Go also provides short variable declaration using the
\texttt{:=} operator:

\begin{lstlisting}
name := "Alice"
score := 95
\end{lstlisting}

The \texttt{:=} operator declares and initialises a new variable, with Go
automatically determining its data type. Short declarations are commonly used
within functions, whereas \texttt{var} is required when a variable needs to be
declared without an initial value or in certain package-level declarations.

Variables can have their values changed after declaration, provided that the
new value is compatible with their data type:

\begin{lstlisting}
score := 95
score = 100
\end{lstlisting}

Go is \textbf{statically typed}, meaning that the type of a variable is
determined at compile time and cannot normally be changed to an unrelated
type during execution. For example:

\begin{lstlisting}
age := 25
age = "twenty-five"    // Compilation error
\end{lstlisting}

This produces a compilation error because \texttt{age} was initially declared
as an integer and cannot subsequently be assigned a string.


\section{Data Types}

A data type defines the kind of data that a variable can contain. Go provides
several built-in data types, including numeric types, strings and Boolean
values. Choosing an appropriate data type helps the compiler determine how
values should be stored and what operations can be performed on them.


\subsection{Integer Types}

The \texttt{int} type is used for whole numbers without decimal places:

\begin{lstlisting}
var age int = 25
\end{lstlisting}

Go also provides explicitly sized integer types, such as \texttt{int8},
\texttt{int16}, \texttt{int32} and \texttt{int64}. The number indicates the
number of bits used to represent the integer.

Unsigned integer types, such as \texttt{uint}, \texttt{uint8} and
\texttt{uint64}, can only represent zero and positive values.


\subsection{Floating-Point Types}

The \texttt{float32} and \texttt{float64} types are used to represent numbers
containing decimal values:

\begin{lstlisting}
var temperature float64 = 21.5
\end{lstlisting}

The \texttt{float64} type provides greater precision than \texttt{float32}
and is commonly used when greater numerical precision is required.


\subsection{Strings}

The \texttt{string} type represents a sequence of characters:

\begin{lstlisting}
var name string = "Alice"
\end{lstlisting}

Strings can also be created using short variable declaration:

\begin{lstlisting}
message := "Hello, world!"
\end{lstlisting}

Strings are commonly used for text such as names, messages and user input.


\subsection{Boolean Values}

The \texttt{bool} type represents a logical value and can contain either
\texttt{true} or \texttt{false}:

\begin{lstlisting}
var loggedIn bool = true
\end{lstlisting}

Boolean variables are commonly used in conditional statements:

\begin{lstlisting}
if loggedIn {
    fmt.Println("User is logged in")
}
\end{lstlisting}


\subsection{Other Go Data Types}

Go also provides more complex data types, including arrays, slices, maps,
structs, pointers, interfaces and functions. These allow programs to
represent collections of values, structured records and more complex
relationships between data.

For example, a slice can store multiple values of the same type:

\begin{lstlisting}
numbers := []int{10, 20, 30}
\end{lstlisting}

A map stores values using keys:

\begin{lstlisting}
ages := map[string]int{
    "Alice": 25,
    "Bob":   30,
}
\end{lstlisting}

Go's strong and static type system allows many type-related errors to be
detected during compilation rather than when the program is running. This
helps improve the reliability and predictability of Go programs.